\subsection{Evaluation}\label{sec:2.3}

% Sources: DETECTION_TASK.md; UCSF_PLAN.md; SYSTEM_PLAN.md; METHODS_PLAN.md; refs varoquaux2018.
\paragraph*{Development and testing without retraining.} Each cohort played one role (Table~\ref{tab:eval}). On MU we
evaluated every model by five repeats of fivefold cross validation grouped by patient and report the mean pooled
held-out AUC, because single splits are unstable in small cohorts \citep{varoquaux2018}. We then fitted each model
once on all 181 grade 4 detection scans of MU, including the 35 without a documented start of radiotherapy, and
applied it unchanged to LUMIERE and UCSF-ALPTDG. At UCSF-ALPTDG every model was planned before any score was compared
with the label. At LUMIERE the combined model was planned and the other scores are exploratory.

% Sources: CLINICIAN_BASELINE_PLAN.md; METHODS_RESULTS.md E4; ref beaulieujones2021.
\paragraph*{The time-and-enhancement model.} We judged each score against a simple model of the time since
radiotherapy and the enhancing tumour volume, which measures what a score adds to evidence a clinician already holds
\citep{beaulieujones2021}. It is a ridge logistic regression on the time window after radiotherapy, the log enhancing
volume and its change since the previous scan, fitted with the same pipeline and folds as PRISM's models and applied
unchanged in the same way. At LUMIERE the window is the rater's flag. The plan fixed the reading in advance. A gain
whose 95\% CI lies above zero counts as a gain, and one whose CI includes zero is reported as no clear gain. To see how
far time alone goes, we also scored scans by timing alone (exploratory, Table~\ref{tab:models}): the days since the
end of radiotherapy and the visit number without any fitted model, and a ridge model on every timestamp input of the
clinical-history model.

% Sources: METHODS_PLAN.md M1; prism/analyses/methods.py; refs kallus2019, janes2008, janes2009.
\paragraph*{Splitting the AUC by time window.} We split each score's AUC into pairs of scans taken at a similar time
and pairs taken at different times. Let $P$ and $N$ be the scans with and without progression, $s$ a score, and $c$
the time window of each scan. The AUC is the share of pairs ranked correctly,
\begin{equation}
\mathrm{AUC} = \frac{1}{|P|\,|N|} \sum_{i \in P} \sum_{j \in N} K(s_i, s_j), \qquad
K(a, b) = \mathbb{1}[a > b] + \tfrac12\, \mathbb{1}[a = b].
\end{equation}
The pairs split into $W = \{(i, j) : c_i = c_j\}$, two scans from the same window, and its complement $B$. With
$A_W$ and $A_B$ the same average over $W$ and over $B$, and $w = |W| / (|P|\,|N|)$,
\begin{equation}
\mathrm{AUC} = w\, A_W + (1 - w)\, A_B
\end{equation}
holds exactly \citep{kallus2019}. $A_W$ is the ranking with the window held fixed, a pair weighted analogue of the
covariate adjusted AUC \citep{janes2008,janes2009}. A score with $A_B$ above $A_W$ ranks pairs across windows better
than pairs within one, which can come from the window or from real differences between the tumours of different
windows. To show how much of $A_B$ the window alone could supply, we also scored every scan by the progression rate of
its window in the training folds (the window only score). At MU the window comes from the treatment dates (during
radiotherapy, within 12 weeks after its end, and later). At LUMIERE it is the rater's flag for scans within three
months of radiotherapy, and at UCSF-ALPTDG it is computed from the start of radiotherapy, assuming a six week course.
Externally, the window only score takes MU's rate for the matching window. The identity was checked on every
bootstrap sample.

% Sources: CHRISTODOULOU_FAITHFUL_PLAN.md (with deviations); CHRISTODOULOU_REPLICATION.md (faithful run);
% ref christodoulou2026.
